\doublespacing % Do not change - required

\chapter{Paper-Based Chapter 1 Title}
\label{ch:paper1}

%%%%%%%%%%%%%%%%%%%%%%%%%%%%%%%%%%%%%%%
% IMPORTANT
\begin{spacing}{1} %THESE FOUR
\minitoc % LINES MUST APPEAR IN
\end{spacing} % EVERY
\thesisspacing % CHAPTER
% COPY THEM IN ANY NEW CHAPTER
%%%%%%%%%%%%%%%%%%%%%%%%%%%%%%%%%%%%%%%

\chapterpublicationnote{This chapter is based on: Author list, ``Paper 1 title,'' venue, year, DOI/link. Replace this note with the full citation, publication status, copyright or permission statement, and a concise author contribution statement.}

\section{Chapter overview}

Do not repeat the paper abstract here. Instead, use this short overview to explain what problem this chapter addresses, how it advances the thesis-level research questions, and how it connects to the surrounding chapters.

\section{Background and motivation}

Summarise only the background needed for this chapter. Add thesis-specific context that may not have appeared in the original paper, especially where it helps the reader understand the wider doctoral project.

\section{Method}

Present the method, system, study design, or analytical approach. Rewrite the paper material into a consistent thesis voice and add details that were omitted from the publication because of page limits.

\section{Experiments and results}

Report the evaluation and findings. Reuse figures or tables only when copyright and permissions allow it, and cite the original publication wherever material is reused.

\section{Discussion}

Discuss what the results show, what they do not show, and how the chapter contributes to the thesis as a whole. This section can be broader than the paper discussion.

\section{Chapter summary and transition}

End with a short synthesis of the chapter's contribution and a transition to the next paper-based chapter.
